\appendix

\section{Формулы для пересчёта на другие модели}
\label{app:formulas}

Обозначим через $d$ число новых токенов в вызове, через $L$ глубину контекста, через $B$ скорость канала в байтах в
секунду.

\textbf{Кэш внимания на токен.} Для $\ell$ слоёв полного внимания с $h$ головами ключей и значений размерности $e$ и
$b$ байтами на число кэш на токен равен $K = 2 \ell h e b$. У Qwen3.8-27B $\ell = 16$, $h = 4$, $e = 256$, а
8-битный формат llama.cpp хранит 34 байта на 32 числа, что даёт 34\,816 байт; замер на сохранённом слоте даёт
34\,844 байта вместе со служебными полями.

\textbf{Дельта.} $D(d) = d \cdot K + S$, где $S$ размер рекуррентного состояния; для моделей со скользящим окном $W$
$D(d) = d \cdot K_f + \min(d, W) \cdot K_w + S$, где $K_f$ и $K_w$ кэш на токен в слоях полного внимания и в слоях окна; у Qwen3.8-27B $S = 156.9$~МБ в 32
битах.

\textbf{Вычисления на токен prefill.} $F(L) \approx 2N + 4 \ell_a L q e$, где $N$ число активных параметров,
$\ell_a$ число слоёв полного внимания, $q$ число голов запросов. Первое слагаемое не зависит от глубины, второе
растёт с ней линейно.

\textbf{Время вызова.} Локальное чтение $t_{\text{loc}} = d / r(L)$, где $r(L)$ замеренная скорость prefill клиента на
глубине $L$. Ход с оффлоадом занимает $t_{\text{off}} = T_0 + t_{\text{up}} + d F(L) / (P \cdot u) + D(d) / B +
t_{\text{ins}} + t_{\text{last}}$, где $T_0$ постоянная задержка вызова, $t_{\text{up}}$ подъём копии сессии на сервере
(абзац о хранении ниже), $P$ пиковая производительность GPU сервера, $u$ доля пика, которую использует движок,
$t_{\text{ins}}$ вставка дельты в кэш клиента, $t_{\text{last}}$ чтение последнего токена порции на клиенте. Клиент
отдаёт порцию серверу, если полное время хода с оффлоадом меньше полного локального: $t_{\text{off}} + t_q <
t_{\text{loc}} + t_{q,\text{loc}}$, где $t_q$ и $t_{q,\text{loc}}$ ожидание в очереди на сервере и на клиенте. Движок
сравнивает именно полные времена; очередей он не моделирует (раздел~\ref{sec:limits}) и берёт оба ожидания
равными нулю.

\textbf{Предел генерации.} Скорость генерации одного потока не выше $W / (A + L K + 2S)$ токенов в секунду, где $W$
полоса памяти клиента, $A$ байты активных весов, читаемых за шаг; рекуррентное состояние за шаг читается и
записывается.

\textbf{Стоимость вызова на сервере.} $c = g \cdot p / u_f + D(d) \cdot c_{\text{tr}}$, где $g$ GPU-секунды на вызов,
$p$ цена GPU-секунды, $u_f$ средняя загрузка парка, $c_{\text{tr}}$ цена байта исходящего трафика.

\textbf{Хранение копии сессии.} Перед чтением сервер поднимает копию в память GPU за $(L K + S) / B$ секунд, где $B$
меньшая из полос PCIe и хранилища; после вызова дописывает новые позиции за $(d + o) K / B$, где $o$ сгенерированные
токены, а состояние перезаписывает при каждом оффлоаде за $S / B$. Если токены вызовов, прочитанных клиентом, сервер
дочитывает при следующем оффлоаде, подъём нужен только вызовам, ушедшим на сервер, и ещё один раз за день на
сессию. Хранение стоит $n \cdot \bar{M} \cdot (a \cdot p_{\text{RAM}} + (1 - a) \cdot p_{\text{NVMe}})$ в месяц, где
$n$ копий на человека, $\bar{M}$ средний размер копии, $a$ доля рабочих часов в месяце, $p$ цена гигабайта в месяц.
В расчётах $n = 4$, $B = 50$~ГБ/с, $p_{\text{RAM}} = \$0.50$ и $p_{\text{NVMe}} = \$0.10$; все четыре величины оценочные.

\section{Осуществимость на прототипе}
\label{app:prototype}

Мы проверили обмен памятью модели на прототипе, построенном на llama.cpp \cite{llamacpp}, с маленькой моделью
Qwen3.5-0.8B той же архитектуры, что Qwen3.8-27B, на процессоре. Связку сервера на CUDA с клиентом на Metal и модель 27B на глубине
сотен тысяч токенов прототип не проверял.
\begin{itemize}
\item Клиент, который вставил дельту и сам посчитал последний токен порции, генерирует те же 64 токена, что и
клиент, прочитавший порцию локально с той же разбивкой на пакеты, на глубинах 512, 2048 и 8192.
\item Против чтения всего промпта одним пакетом распределения следующего токена расходятся не больше чем на 0.00204
ната на токен (среднее расхождение Кульбака--Лейблера при подстановке верных токенов), а самый вероятный токен
совпадает в 63 или 64 позициях из 64. Свободная жадная генерация при другой разбивке на пакеты может разойтись в
отдельной позиции; так ведёт себя и чисто локальный запуск, и к переносу памяти это не относится.
\item Вставка дельты занимает 4--5 миллисекунд и не растёт с глубиной от 512 до 8192 токенов.
\item Пятнадцать видов испорченных дельт клиент отклоняет, и его кэш остаётся побайтно прежним.
\item Журнал сессии сокращает число повторно прочитанных токенов в разговоре из трёх ходов с 1530 до 524.
\end{itemize}
Код прототипа открыт: \url{https://github.com/relux-works/hybrid-inference}.
